% Generated by task tex-libraries from dossiers/lib-others.md, section A (and E.1, E.4, F.3 for
% the new pin). Snippets are copied from the dossier, which copied them from the sources.
\section{VCVio: oracle computations, and what a probability means}\label{sec:gen-lib-vcvio}

VCVio is the library of \emph{oracle computations} on which ArkLib states every security notion.
Every snippet of this section is from VCVio at the old pin \code{f9dc47d9} (\enquote{feat(complexity):
add honest oracle-PPT foundations and measure-native OTP (\#545)}), read under
\file{.lake/packages/VCVio/}; \cref{sec:gen-lib-vcvio-new} says what the new pin \code{a4232d08}
changes.

\subsection{Oracles and computations}\label{sec:gen-lib-vcvio-comp}

\begin{define}{\lean{OracleSpec}: the list of oracles a computation may call}
An oracle specification names a set of oracles by an index type \lean{ι} and gives, for each
index (each possible query), the type of the answer. It is nothing more than a function to
types. \lean{A →ₒ B} is the specification of one oracle with queries in \lean{A} and answers in
\lean{B}; \lean{[]ₒ} is the empty specification (no oracle at all); \lean{unifSpec} is the source
of all randomness, the oracle of uniform choices from \lean{Fin (n+1)}.
\end{define}
\begin{leancode}
/-- An `OracleSpec ι` specifies a set of oracles indexed by `ι`.
Defined as a map from each input to the type of the oracle's output. -/
def OracleSpec (ι : Type u) : Type (max u (v + 1)) :=
  ι → Type v
\end{leancode}
\src{\file{VCVio/OracleComp/OracleSpec.lean:27-30}, VCVio \code{f9dc47d9}}

\begin{leancode}
/-- Specifies access to no oracles, using the empty type as the indexing type. -/
@[reducible] def emptySpec : OracleSpec PEmpty := PEmpty →ₒ PEmpty
notation "[]ₒ" => emptySpec
\end{leancode}
\src{\file{VCVio/OracleComp/OracleSpec.lean:215-217}, VCVio \code{f9dc47d9}}

\begin{leancode}
/-- Access to oracles for uniformly selecting from `Fin (n + 1)` for arbitrary `n : ℕ`.
By adding `1` to the index we avoid selection from the empty type `Fin 0 ≃ empty`. -/
@[inline, reducible] def unifSpec : OracleSpec ℕ :=
  OracleSpec.ofFn fun n => Fin (n + 1)
\end{leancode}
\src{\file{VCVio/OracleComp/OracleSpec.lean:233-236}, VCVio \code{f9dc47d9}}

(The notation is \code{notation:25 (name := singletonSpec) A:25 " →ₒ " B:26 =>
OracleSpec.ofFn (ι := A) (fun _ => B)}, \file{OracleSpec.lean:84-85}.)

\begin{define}{\lean{OracleComp} and \lean{ProbComp}: programs that call oracles}
\lean{OracleComp spec α} is a program returning an \lean{α} that may call the oracles of
\lean{spec}. It is a \emph{free monad}: a syntax tree whose nodes are \enquote{query $t$, then
continue with the answer}. It has no meaning by itself; a meaning is supplied by interpreting
the queries (\cref{sec:gen-lib-vcvio-dist}). A \emph{probabilistic} computation,
\lean{ProbComp}, is one whose only oracle is the uniform-choice oracle, with the primitive coin
\lean{$[0..n]} (\code{def uniformFin (n : ℕ) : ProbComp (Fin (n + 1)) := unifSpec.query n},
\file{ProbComp.lean:55-58}).
\end{define}
\begin{leancode}
/-- `OracleComp spec α` represents computations with oracle access to oracles in `spec`,
where the final return value has type `α`, represented as a free monad over the `PFunctor`
corresponding to `spec.` -/
@[reducible]
def OracleComp {ι : Type u} (spec : OracleSpec.{u, v} ι) :
    Type w → Type (max u v w) :=
  PFunctor.FreeM spec.toPFunctor
\end{leancode}
\src{\file{VCVio/OracleComp/OracleComp.lean:23-29}, VCVio \code{f9dc47d9}}

\begin{leancode}
/-- Simplified notation for computations with no oracles besides random inputs.
This specific case can be used with `#eval` to run a random program, see `OracleComp.runIO`.
NOTE: Need to decide if this should be more opaque than `abbrev`, seems like no as of now.. -/
abbrev ProbComp : Type → Type := OracleComp unifSpec
\end{leancode}
\src{\file{VCVio/OracleComp/ProbComp.lean:42-45}, VCVio \code{f9dc47d9}}

\paragraph{Failure.} An \lean{OracleComp} tree has no failure node: it always returns. Failure (a
verifier rejecting, a guard not met) is added by Lean's \lean{OptionT} monad transformer,
\lean{OptionT m α = m (Option α)}, where \lean{none} means \enquote{failed}. ArkLib's verifiers
run in \lean{OptionT (OracleComp oSpec)}. VCVio sends \lean{none} to the \enquote{missing mass}
of a sub-distribution (\lean{PMF} is Mathlib's type of probability mass functions: a function
\lean{α → ℝ≥0∞} summing to 1):
\begin{leancode}
/-- A subprobability mass function is a function `α → ℝ≥0∞` such that values have an infinite
sum at most `1` represented by applying an `OptionT` transformer to the `PMF` monad.
The new `failure`/`none` value holds the "missing" mass to reach the total sum of `1`. -/
def SPMF : Type u → Type u := OptionT PMF
\end{leancode}
\src{\file{ToMathlib/ProbabilityTheory/SPMF.lean:55-58}, VCVio \code{f9dc47d9}}

\begin{leancode}
/-- Lift a `MonadLiftT m SPMF` instance to `MonadLiftT (OptionT m) SPMF`. Failure in `OptionT`
contributes to the failure mass of the resulting `SPMF`. -/
noncomputable instance instMonadLiftTSPMF (m : Type u → Type v) [Monad m]
    [MonadLiftT m SPMF] [LawfulMonadLiftT m SPMF] :
    MonadLiftT (OptionT m) SPMF where
  monadLift x := OptionT.mapM' (MonadHom.ofLift m SPMF) x
\end{leancode}
\src{\file{VCVio/EvalDist/Instances/OptionT.lean:126-131}, VCVio \code{f9dc47d9}}

\begin{define}{\lean{QueryImpl} and \lean{simulateQ}: answering the queries}
A \lean{QueryImpl spec m} answers every query of \lean{spec} by a computation in another monad
\lean{m}; \lean{simulateQ impl oa} replaces every query node of \lean{oa} by its answer. This is
how ArkLib's challenger answers \enquote{give me challenge $i$} by a uniform sample, and how an
oracle is implemented from the message it stands for.
\end{define}
\begin{leancode}
/-- A monadic handler for the polynomial interface induced by `spec`.

Concretely, this maps every oracle input `x` to a computation returning an
answer of type `spec.Range x`. It extends first to `OracleQuery spec` by
applying the continuation, then to `OracleComp spec` by preserving `pure` and
`bind`; see `QueryImpl.mapQuery` and `simulateQ`. -/
@[reducible] def QueryImpl {ι} (spec : OracleSpec ι) (m : Type u → Type v) :=
  (x : spec.Domain) → m (spec.Range x)
\end{leancode}
\src{\file{VCVio/OracleComp/SimSemantics/QueryImpl/Basic.lean:30-37}, VCVio \code{f9dc47d9}}

\begin{leancode}
/-- Given an implementation of `spec` in the monad `r`, convert an `OracleComp spec α` to a
implementation in `r α` by substituting `impl t` for `query t` throughout. -/
def simulateQ {ι} {spec : OracleSpec ι} {r : Type u → Type _} [Monad r]
    (impl : QueryImpl spec r) {α : Type u} (mx : OracleComp spec α) : r α :=
  PFunctor.FreeM.liftM impl mx
\end{leancode}
\src{\file{VCVio/OracleComp/SimSemantics/SimulateQ.lean:27-31}, VCVio \code{f9dc47d9}}

\subsection{The distribution of a computation, \texorpdfstring{\lean{Pr[…]}}{Pr[...]}, and the support}\label{sec:gen-lib-vcvio-dist}

\begin{define}{The distribution of a computation}
The distribution of an \lean{OracleComp} is obtained by \lean{simulateQ} with every query
answered by a fixed distribution on its answers. For \lean{unifSpec} (hence every
\lean{ProbComp}) that distribution is the uniform distribution on \lean{Fin (n+1)}, declared
once. The probabilities are then read off that sub-distribution.
\end{define}
\begin{leancode}
@[reducible] noncomputable def IsUniformSpec.ofFintypeInhabited
    {ι : Type u} (spec : OracleSpec ι)
    [hF : spec.Fintype] [hI : spec.Inhabited] : IsUniformSpec spec where
  toPMF t := PMF.uniformOfFintype (spec.Range t)
  fintype := hF
  inhabited := hI
  toPMF_eq_uniform _ := rfl

noncomputable instance : IsUniformSpec unifSpec := IsUniformSpec.ofFintypeInhabited _
\end{leancode}
\src{\file{VCVio/OracleComp/EvalDist.lean:71-80}, VCVio \code{f9dc47d9}}

\begin{leancode}
lemma evalSPMF_eq_simulateQ [IsProbabilitySpec spec] (mx : OracleComp spec α) :
    𝒮[mx] = simulateQ IsProbabilitySpec.toPMF mx := rfl
\end{leancode}
\src{\file{VCVio/OracleComp/EvalDist.lean:146-147}, VCVio \code{f9dc47d9}}

\begin{leancode}
def probOutput [MonadLiftT m SPMF] (mx : m α) (x : α) : ℝ≥0∞ :=
  evalSPMF mx x
...
noncomputable def probEvent [MonadLiftT m SPMF] (mx : m α) (p : α → Prop) : ℝ≥0∞ :=
  (evalSPMF mx).run.toOuterMeasure (some '' {x | p x})

/-- Probability that a computation `mx` will fail to return a value. -/
def probFailure [MonadLiftT m SPMF] (mx : m α) : ℝ≥0∞ :=
  (evalSPMF mx).run none

/-- Probability that a computation returns a particular output. -/
notation "Pr[= " x " | " mx "]" => probOutput mx x

/-- Probability that a computation returns a value satisfying a predicate. -/
macro (name := probEventNotation) "Pr[ " p:term " | " mx:term "]" : term =>
  `(probEvent $mx $p)

/-- Probability that a computation fails to return a value. -/
notation "Pr[⊥" " | " mx "]" => probFailure mx
\end{leancode}
\src{\file{VCVio/EvalDist/Defs/Basic.lean:114-140}, VCVio \code{f9dc47d9}}

So \lean{Pr[= x | mx]} is the probability that \lean{mx} returns \lean{x}, \lean{Pr[p | mx]}
the probability that it returns a value satisfying \lean{p} (failure never satisfies an event),
\lean{Pr[⊥ | mx]} the probability that it fails; all are extended non-negative reals
(\lean{ℝ≥0∞}). The \textbf{support} is the set of values a computation can return, computed
syntactically (every answer allowed at every query):
\begin{leancode}
/-- The set of possible outputs of running the monadic computation `mx`. -/
def support [MonadLiftT m SetM] {α : Type u} (mx : m α) : Set α :=
  SetM.run (liftM mx)
\end{leancode}
\src{\file{VCVio/EvalDist/Defs/Support.lean:42-44}, VCVio \code{f9dc47d9}}

\subsection{Uniform sampling, with uniformity as a law}\label{sec:gen-lib-vcvio-sample}

\begin{define}{\lean{SampleableType} and \lean{$ᵗ}: uniform sampling}
A \lean{SampleableType β} instance is a sampler \lean{selectElem : ProbComp β} \emph{and two
proofs}: every element is reachable, and any two elements have equal probability.
\lean{$ᵗ β} is that sampler. Uniformity is part of the structure: because a \lean{ProbComp}
never fails, the probabilities sum to one, and VCVio derives the uniform law for every instance.
Consequently \textbf{no instance of \lean{SampleableType} can be non-uniform}, and two instances
on the same finite type have the same distribution (proved in the probe \code{SamplerDiamond},
\lean{sampler_unique}; \cref{sec:gen-l0-samplers}).
\end{define}
\begin{leancode}
/-- A `SampleableType β` instance means that `β` is a finite inhabited type,
with a computation `selectElem` that selects uniformly at random from the type.
...
class SampleableType (β : Type) where
  selectElem : ProbComp β
  mem_support_selectElem (x : β) : x ∈ support selectElem
  probOutput_selectElem_eq (x y : β) : Pr[= x | selectElem] = Pr[= y | selectElem]

/-- Select uniformly from the type `β` using a type-class provided definition.
NOTE: naming is somewhat strange now that `Fintype` isn't explicitly required. -/
def uniformSample (β : Type) [h : SampleableType β] : ProbComp β := h.selectElem

notation:90 "$ᵗ " α:91 => uniformSample α
\end{leancode}
\src{\file{VCVio/OracleComp/Constructions/SampleableType.lean:37-53}, VCVio \code{f9dc47d9}}

\begin{leancode}
/-- Every element of a uniform sample over a `Fintype` has output probability `card⁻¹`. -/
@[simp, grind =]
lemma probOutput_uniformSample [Fintype α] (x : α) :
    Pr[= x | $ᵗ α] = (Fintype.card α : ℝ≥0∞)⁻¹ := by
\end{leancode}
\src{\file{SampleableType.lean:57-66}, VCVio \code{f9dc47d9}}

\begin{leancode}
@[simp, grind =]
lemma probEvent_uniformSample [Fintype α] (p : α → Prop) [DecidablePred p] :
    Pr[ p | $ᵗ α] = (Finset.univ.filter p).card / Fintype.card α := by
\end{leancode}
\src{\file{SampleableType.lean:224-228}, VCVio \code{f9dc47d9}}

The constructors leanerVM uses are the base case, a uniform \lean{Fin n}, which is the
primitive coin; the transport along a bijection, which re-proves the laws; vectors, with
independent coordinates (laws proved at \file{:317-336}); and any \lean{FinEnum} type, that is,
a type with a computable enumeration \lean{α ≃ Fin n} (this instance matters for the diamond on
$\K$, \cref{sec:gen-l0-samplers}).
\begin{leancode}
@[reducible] def SampleableType.Fin (n : ℕ) : SampleableType (Fin (n + 1)) where
  selectElem := $[0..n]
  mem_support_selectElem := by simp
  probOutput_selectElem_eq := by simp

instance (n : ℕ) [hn : NeZero n] : SampleableType (Fin n) :=
  match n, hn with
  | _ + 1, _ => SampleableType.Fin _
\end{leancode}
\src{\file{SampleableType.lean:232-239}, VCVio \code{f9dc47d9}}

\begin{leancode}
/-- A type equivalent to a `SampleableType` is also `SampleableType`. -/
@[reducible] def SampleableType.ofEquiv {α β : Type} [SampleableType α] (e : α ≃ β) :
    SampleableType β where
  selectElem := e <$> ($ᵗ α)
  mem_support_selectElem x := by simp
  probOutput_selectElem_eq x y := by
    rw [probOutput_map_equiv, probOutput_map_equiv]
    exact probOutput_uniformSample_inj α (e.symm x) (e.symm y)
\end{leancode}
\src{\file{SampleableType.lean:260-267}, VCVio \code{f9dc47d9}}

\begin{leancode}
/-- Select a uniform element from `Vector α n` by independently selecting `α` at each index. -/
instance (α : Type) (n : ℕ) [SampleableType α] : SampleableType (Vector α n) where
  selectElem := by induction n with
  | zero => exact pure #v[]
  | succ m ih => exact Vector.push <$> ih <*> ($ᵗ α)
\end{leancode}
\src{\file{SampleableType.lean:313-316}, VCVio \code{f9dc47d9}}

\begin{leancode}
instance FinEnum.SampleableType (α : Type)
    [h : FinEnum α] [Nonempty α] : SampleableType α := by
  have : NeZero (FinEnum.card α) := NeZero.mk FinEnum.card_ne_zero
  exact SampleableType.ofEquiv h.equiv.symm
\end{leancode}
\src{\file{SampleableType.lean:269-273}, VCVio \code{f9dc47d9}}

ArkLib's challenger answers the challenge query of round $i$ by exactly this sampler:
\begin{leancode}
def challengeQueryImpl {pSpec : ProtocolSpec n} [∀ i, SampleableType (pSpec.Challenge i)] :
    QueryImpl ([pSpec.Challenge]ₒ'challengeOracleInterface) ProbComp :=
  fun q => $ᵗ (pSpec.Challenge q.1)
\end{leancode}
\src{\file{ArkLib/OracleReduction/ProtocolSpec/Basic.lean:732-734}, ArkLib \code{dca90385}}

\subsection{What a reader trusts when a theorem is a VCVio probability bound}\label{sec:gen-lib-vcvio-trust}

A statement \lean{Pr[bad | game] ≤ ε} in VCVio unfolds to: interpret the program \lean{game} (a
free-monad tree) by answering each uniform-choice query $n$ with Mathlib's
\lean{PMF.uniformOfFintype (Fin (n+1))} and each \lean{OptionT} failure by the missing mass,
compose these distributions by Mathlib's \lean{PMF} bind, and measure the event. The reader
trusts:
\begin{enumerate}
\item that the \textbf{program} \lean{game} is the intended experiment (for the proof system,
  that ArkLib's round-by-round knowledge-soundness game is the right game; that is the subject of
  \cref{sec:gen-lib-ark-sec});
\item the \textbf{three definitions} above (\lean{IsUniformSpec unifSpec} answers with the
  uniform distribution; \lean{evalSPMF = simulateQ toPMF}; \lean{probEvent} measures
  \lean{some '' {x | p x}}), which are short and standard; and Mathlib's \lean{PMF} (bind,
  \lean{uniformOfFintype}, \lean{toOuterMeasure});
\item that each \textbf{challenge type's} \lean{SampleableType} instance is the intended
  sampler. Since any instance is uniform, what remains to trust is only that the challenge
  \emph{type} is right ($\E$, not a subtype or a smaller field).
  \lean{Fintype.card E = 2^192} (\lean{card_E}, \cref{sec:gen-l0-catalogue}) fixes it.
\end{enumerate}
Nothing else in VCVio is on the trusted path of a statement: its many lemmas are proofs, checked
by the kernel. (VCVio at the old pin also has a measure-theoretic semantics \lean{𝒟[…]}; the
\lean{Pr[…]} notation used by ArkLib and leanerVM is the discrete \lean{SPMF} one quoted above,
\file{Defs/Basic.lean:110-125}.)

\subsection{At the new pin \texorpdfstring{\code{a4232d08}}{a4232d08}}\label{sec:gen-lib-vcvio-new}

\begin{sloppypar}
At VCVio \code{a4232d08} the primary semantics is a Mathlib \lean{Measure} (\lean{𝒟[mx]}), a
fold of the free monad with per-query measures, proved to agree with the old mass-function
semantics on discrete answer spaces (the agreement is proved in
\file{OracleComp/EvalDist/UniformCompatibility.lean:47} and in
\file{EvalDist/PFunctorMeasure.lean:96-106}). \lean{probOutput}, \lean{probEvent} and
\lean{probFailure} are deprecated since 2026-09-13 (\file{VCVio/EvalDist/Defs/Basic.lean:90-110}
at \code{a4232d08}, \code{@[deprecated "VCVio retiring probability API: use 𝒟[mx] {x | p x}"]}) in favour of
\lean{Pr{let x ← c}[P x]}, a macro for the measure of the event \lean{{True}} under the
computation returning the proposition:
\end{sloppypar}
\begin{leancode}
macro_rules (kind := prEvent)
  -- `doSeqBracketed`
  | `(Pr{{$items*}}[$t]) => `(𝒟[do $items:doSeqItem* return $t:term] {True})
  -- `doSeqIndent`
  | `(Pr{$items*}[$t]) => `(𝒟[do $items:doSeqItem* return $t:term] {True})
\end{leancode}
\src{\file{VCVio/EvalDist/ProbabilityNotation.lean:65-69}, VCVio \code{a4232d08}}

\lean{SampleableType}'s law is now stated directly as \enquote{the output measure is uniform},
and sampler irrelevance is proved upstream
(\lean{SampleableType.prEvent_uniformSample_inst_irrel}, \file{NativeMeasure.lean:280-284}):
\begin{leancode}
class SampleableType (β : Type) where
  /-- The canonical sampler of `β`, written `$ᵗ β`. -/
  selectElem : ProbComp β
  /-- The canonical sampler denotes the uniform measure on `β`. -/
  evalDist_selectElem_eq_uniform :
    ∀ [MeasurableSpace β] [MeasurableSingletonClass β],
      𝒟[selectElem] = uniformOn Set.univ
\end{leancode}
\src{\file{VCVio/OracleComp/Constructions/SampleableType/Basic.lean:51-57}, VCVio \code{a4232d08}}

\begin{sloppypar}
What a reader trusts changes accordingly: the measure fold
(\path{VCVio/EvalDist/PFunctorMeasure/Core.lean:226-229}, \lean{instEvalDistSemanticsFreeM}; a
bind whose continuation is not almost-everywhere measurable denotes $0$, \file{:102}, which
cannot arise for discrete answer types) plus the compatibility proof. leanerVM at \code{144c5aa}
restated its two probability bridges in the new notation; for the counting bounds the two
statements are the same mathematical statement (\cref{sec:gen-l0-bounds}).
\end{sloppypar}
